\section{The same explicit cutoff for arbitrary linear characters}
\label{sec:linear-nonweil-cutoff}

\paragraph{Status.}
New derivation, completely read and checked by the author, root and
an independent peer, with separately executed exact constant certificates. The published inputs are the same two articles as in
D-UNW-01, with the full general-linear dual-pair proof specified below.
No new character table or representation classification is assumed.

\begin{theorem}[D-LNW-01: an explicit linear normal-core rank cutoff]
\label{D-LNW-01}
Set $N_0=10^{61}$. If $n\ge N_0$, $q\ge2$ and
$\chi\in\operatorname{Irr}(\operatorname{SL}_n(q))$ has
character level at least two and degree $a$, then
\[
            |\chi(g)|\le\frac4{15}a
           \quad(g\notin Z(\operatorname{SL}_n(q))),
           \qquad a>q^n>30.                                \tag{LNW.1}
\]
Every finite-order nonscalar normalizer of its represented image has
projective support at least $11a/60$.

An irreducibly acting normal central-product core all of whose factors
have such large-natural-dimension non-Weil linear representations
cannot supply period rank nine in the sixfold finite-cover problem.
Arbitrary mixtures among these factors, scalar and outer intertwiners,
central quotients and compatible factor permutations are allowed.

Consequently, in the transitive primitive factorwise
tensor-indecomposable normal-core situation of D-BI-19, a linear
simple quotient $\operatorname{PSL}_n(q)$ must have
\[
                            n<10^{61}.                      \tag{LNW.2}
\]
The level-one alternative is excluded by the already proved linear
Weil theorem D-BI-23. In the nonsplit main-source scope of D-BI-20
the remaining general linear gate also has $q\le27$.
\end{theorem}

\begin{proof}
Write $j$ for the character level. The definition and its independence
are Definition~2 and Lemma~7.1 of \cite[pp.~51--52]{UNW:GLT2020}.
The uniform support estimate of \cite[Theorem~5.5, accepted-author
pp.~25--26]{UNW:LT2024} applies to $\operatorname{SL}_n(q)$ as
well as the unitary group, with the same certified constant
$\sigma_0=10^{-29}$. Theorem~7.3(i)--(ii) of
\cite[p.~52]{UNW:GLT2020} gives
\[
 a\ge\frac{q^{j(n-j)}}{q-1},\qquad
 j\ge n/2\ \Longrightarrow\
 a>\frac{q^{n^2/4-2}}{q-1}.                                \tag{LNW.3}
\]
These imply the weaker common bounds (UNW.4) and hence (UNW.8).
Put $J_0=10^{30}$. If $j>J_0$, the complete argument in
(UNW.9) applies verbatim and gives $|\chi(g)|/a<1/4$.
This does not require a same-degree general-linear extension in
the high-level range.

For $2\le j\le J_0$, we have $j<n/2$. Corollary~8.6, including
its full proof on pp.~67--68 of \cite{UNW:GLT2020}, gives an
irreducible extension to $\operatorname{GL}_n(q)$ of the same
degree; determinant twisting selects true level $j$ and is trivial
on the special linear group. The general-linear dual-pair argument
in the complete \S8.1, pp.~55--61, identifies this extension with
$D_\alpha^\circ$ for an irreducible character $\alpha$ of
$Q=\operatorname{GL}_j(q)$. The condition in Theorem~1.1(ii)
is automatic for $j\le n/2$.

Here is the precise stronger estimate used, rather than the
numerical $1.76$ corollary. Let
$R_j=\sqrt{8q^{j^2+j}}$, and let $k\le n-1$ be the largest
natural eigenspace dimension of a noncentral $g$. Equation (8.11),
p.~60, gives
\[
\begin{split}
 a&\ge\frac{\alpha(1)}{|Q|}
  \left(q^{nj}-|Q|q^{n(j-1)}-|Q|R_jq^{n(j-1)}\right),\\
 |\chi(g)|&\le\frac{\alpha(1)}{|Q|}
  \left(q^{kj}+|Q|q^{(n-1)(j-1)+1}
                         +|Q|R_jq^{n(j-1)}\right).
\end{split}                                                    \tag{LNW.4}
\]
The complete derivation of the tail bound is (8.9)--(8.10),
pp.~59--60. The earlier geometric dual-pair argument proves the
irreducibility and exhaustion assertion; no nonnegative virtual
character is silently substituted for an irreducible one.

The elementary order formula gives $|Q|\le q^{j^2}$, and
$R_j\le3q^{j^2}$ because $8<9$ and $j^2+j\le2j^2$ for $j\ge1$.
Since $N_0\ge2J_0^2+9$, put
\[
                   u=q^{2j^2-n}\le2^{-9},\qquad d=4u.
\]
After division by $\alpha(1)q^{nj}/|Q|$, the lower bound in
(LNW.4) is at least $1-d$, since
\[
 |Q|(1+R_j)q^{-n}
 \le q^{j^2-n}+3q^{2j^2-n}\le4u.
\]
The normalized upper bound is at most $q^{-j}+d$: the middle
term has exponent at most $j^2-n-j+2\le j^2-n$.
Consequently
\[
 \frac{|\chi(g)|}{a}
 \le\frac{1/4+1/128}{1-1/128}
 =\frac{33}{127}<\frac{26}{99}<\frac4{15}.                 \tag{LNW.5}
\]
This includes $q=2,j=2$. The degree separation $a>q^n$ for
levels at least two follows from the same inequalities at the end
of D-UNW-02, using (LNW.3). Level-one characters have degree at
most $q^n$ by Theorem~7.3(i); they are the linear Weil characters,
as explicitly recorded on p.~61 of \cite{UNW:GLT2020}.

For a nonscalar normalizer $T$, choose a core element $x$ with
$[T,x]\ne1$ by irreducibility. A nonidentity scalar commutator
has ordinary moving rank $a$. Otherwise (LNW.1) and the unitary
eigenvalue trace inequality give ordinary moving rank at least
$(1-4/15)a/2$. In either case
$\operatorname{rank}([T,x]-1)\le2\nu(T)$, so
$\nu(T)\ge11a/60$. This treats every compatible outer
automorphism directly, without imposing a field or graph
stabilizer hypothesis.

For tensor products, normalized eigenspace multiplicities are
convolutions, and a nontrivial factor-permutation cycle has at
least one-quarter projective support by the complete weighted
word-orbit argument in D-UNW-03. The common bound is therefore
$11e/60$ on total degree $e$. The same actual intermediate-base
Euler identity (UNW.15) forces genus zero and at least twelve
genuine inertia points for period rank nine, and then
$2e+6\ge12(11e/60)$ gives $e\le30$. Each factor already has
degree greater than thirty, a contradiction.

At natural dimension $n\ge N_0$, all quasisimple groups with
simple quotient $\operatorname{PSL}_n(q)$ are central quotients
of $\operatorname{SL}_n(q)$; the complete exceptional-cover list
in \cite[Chapter~5, Table~1, p.~35]{BI:HHM} has no such parameter.
Thus all factor characters lift to this ordinary group. In the
transitive-core case their degrees are equal. Degree separation
forces all of them to have level one or all to have level at least
two. The latter case has just been excluded; the former is exactly
the range of D-BI-23, with no large-rank exception. This proves
(LNW.2). The field cutoff invokes D-BI-20 only with its original
nonsplit and primitive-core hypotheses.
\end{proof}

\paragraph{Verification boundary.}
The entire cited general-linear \S8.1, including Lemma~8.2 and
the exact estimate (8.11), was read from the complete published PDF.
The standalone \texttt{check\_linear\_constants.py} executes the
new exact low-level estimates and checks the reduction to the already
certified high-level constants. Neither this corollary nor the
certificate enumerates the finite remaining range, claims that a
remaining profile is sufficient, or admits additional uncovered
small factors in its mixed-core statement.
